\scanpage{198}
\mspage{L 17}
% archive sheet number: 197

entrevu seulement et encore moins vu et touché, pour peu à peu se \deleted{dégager de son manteau de brumes, et se} \deleted{\inserted{d'ombre et}} décanter de l'impondérable, se dégager de son manteau d'ombre et de brumes et prendre forme et chair. \deleted{Ce sont} \inserted{C'est} \deleted{ces stades} \inserted{cette partie}-là du travail, de piètre apparence pour ne pas dire foireux, qui \deleted{\unsure{en représentent aussi les plu|en représentent aussi les plus|à y regarder de plus près}} \inserted{en \deleted{sont} est aussi la plus} névralgique - celle où, véritablement quelque chose de \mbox{n\,o\,u\,v\,e\,a\,u} fait son apparition, par l'effet d'une attention intense, d'une sollicitude, d'un respect pour cette chose fragile, infiniment délicate, sur le point de naître. C'est la partie créatrice entre toutes - celle\deleted{\unsure{s}} de la conception et de la gestation d'un \unsure{fragile|gragile} embryon dans \unsure{hes|les} chaudes ténèbres de la matrice nourricière, \deleted{élaborant en silence} depuis l'invisible double gamète \deleted{initiale} originelle devenant embryon et se transformant au fil des jours et des mois, par un travail obscur et intense, \unsure{irvisible|invisible} et sans apparence, en un nouvel être en chair et en os.

\unsure{C'est|Cest} là aussi la partie \inserted{éminemment} ``yin'', ``féminine\deleted{\unsure{,}}'' du travail de découverte - alors que \deleted{la partie} l'aspect complémentaire ``yang'', ou masculin, s'apparenterait plutôt \deleted{\unsure{aux|auxx}} au travail à coups de marteaux ou de masse, sur un burin bien affûté ou sur un coin de bon acier trempé. L'un et l'autre aspect a sa raison d'être et sa fonction, en symbiose inséparable \inserted{l'un} avec l'autre - ou pour mieux dire, ce sont là \uline{l'épouse} et \uline{l'époux} du couple des deux forces cosmiques originelles, dont l'étreinte \deleted{\inserted{sans cesse} renouvelée fait} sans cesse renouvelée fait resurgir sans cesse les obscurs labeurs créateurs de la conception, de la gestation et de la naissance - de la naissance de \mbox{l'\,e\,n\,f\,a\,n\,t} , de la chose nouvelle.

La deuxième chose sur laquelle je sentais le besoin de m'exprimer, dans ma fameuse ``introduction'' \deleted{philosophique et} personelle et ``philosophique'' à un texte mathématique, c'était au sujet de la nature du travail créateur justement. Je me rendais compte, depuis des années, que cette nature était généralement \unsure{ignorée|ngnorée}, occultée par des clichés à tout venant et par des peurs ancestrales. A quel point il en est bien ainsi, je l'ai découvert \unsure{progressivement|progressigement}, \unsure{au|an} fil des jours et des mois, tout au cours de la réflexion et de ``l'enquêté'' poursuivie dans Récoltes et Semailles. C'est dès les tout débuts de cette réflexion, au mois de juin 1983, que je suis saisi pour la \unsure{première|premmère} fois (pour l'oublier \inserted{un peu} par la suite ...) par la portée de ce fait d'anodine apparence, et pourtant asses stupéfiant, pour peu seulement qu'on s'y arrête tant soit peu: que cette partie ``créatrice entre toutes'' dont je viens de parler dans le travail de découverte, \uline{ne transparaît pratiquement nulle part} dans les textes qui sont censés
